1. Introduction




Ecuador has experienced an unprecedented explosion of lethal violence. Intentional homicides recorded by the national police registry rose from 996 in 2018 to 9,283 in 2025 -- a nine-fold increase and a national rate of approximately 51 per 100,000 inhabitants -- with 4,154 additional deaths in the first half of 2026 alone (DINASED, Ministry of the Interior). The escalation accelerated abruptly from 2,495 homicides in 2021 to 8,248 in 2023, plateauing at historically high levels in 2024--2025. The surge coincided with Ecuador's consolidation as a transit and storage hub for cocaine shipments through its Pacific ports, the territorial expansion of transnational organized crime, and a sequence of executive responses culminating in the declaration of an "internal armed conflict" (Decrees 111, 9 January 2024; 218, 7 April 2024; 55, 16 July 2025; and 424, 18 June 2026). Public debate has frequently attributed the wave to social causes -- particularly the deterioration of the labor market -- while the government has framed the crisis primarily as a war against organized crime.


This paper subjects the labor market channel to systematic empirical scrutiny at the provincial level, the finest geographic scale for which official data are publicly available. Ecuador's labor market displays a paradoxical profile: open unemployment is low (3.48\% in ENEMDU 2024-II), but underemployment is high (20.97\%) and adequate employment covers only 35.01\% of the labor force -- a combination typical of economies with large informal sectors, where the relevant margin for criminal opportunity cost is the quality and stability of licit income rather than the unemployment rate per se. The central question is twofold. First, how strongly are provincial labor market conditions associated with provincial homicide rates between 2018 and 2026, once geographic and temporal heterogeneity are taken into account? Second, is there evidence that the state's exceptional response (Decrees 111 and successors) altered the trajectory of homicides?


We address these questions with a province--quarter panel of 24 provinces (2018Q1--2026Q2, 816 observations) constructed exclusively from official sources: incident-level homicide records from the DINASED registry of the Ministry of the Interior, population projections from INEC, provincial labor indicators computed from ENEMDU microdata (2024Q2, the only provincial cross-section currently available), and CONALI/IGM provincial boundaries. The empirical strategy combines descriptive cross-sectional correlations (N = 24), regression specifications with provincial and period fixed effects and their identification limits, a dynamic event study around Decreto 111 with a placebo event in 2021Q1, and a synthetic control method (SCM) for Esmeraldas. We also document the data extensions required for credible causal inference.


The paper makes four contributions. First, it provides the first systematic, reproducible description of the homicide wave at the provincial level using official administrative data, including an external validation against the independent bulletin of the Observatorio Ecuatoriano de Crimen Organizado (OECO). Second, it delivers a transparent and falsifiable assessment of the labor market channel, showing that the association between labor indicators and homicides is weak in cross-section and statistically indistinguishable from zero in panel regressions -- a pattern coherent with the international literature, which finds unemployment effects on property crime but not on homicides (Chiricos, 1987; Raphael \& Winter-Ebmer, 2001; Edmark, 2005; Draca \& Machin, 2015). Third, it applies modern quasi-experimental tools -- event studies with placebo designs and synthetic controls -- to the timing of the security response, showing that the wave preceded the decrees and that no significant differential change followed them, while being explicit about the inferential limits of these tools with 24 clusters. Fourth, it translates the findings into concrete, evidence-labeled recommendations for senior policymakers.


Three caveats frame the analysis. First, the labor panel covers a single ENEMDU cross-section (2024-II): provincial labor indicators are time-invariant in our data, which prevents the estimation of long-run elasticities and limits identification to cross-sectional variation. Second, this is not a causal exercise: fixed effects do not eliminate omitted variables such as police deployment, armed group presence, migration, or drug market shocks, and reverse causality (crime destroying employment) is a documented concern (IMF, 2024). Third, administrative homicide records are "subject to variations" according to the official metadata itself. The paper reports these limitations explicitly and repeatedly, because the policy stakes of misreading correlation as causation are high.


The remainder of the paper is organized as follows. Section 2 describes the data sources, the construction of the panel, and the descriptive evidence. Section 3 presents the econometric framework. Section 4 reports the results. Section 5 discusses the findings in light of the verified international literature and derives policy implications with explicit evidence labels. Section 6 concludes.